% \begin{document}
\section{Large-step Operational Semantics}

In the last lecture we defined a semantics for our language of arithmetic expressions using a
small-step evaluation relation $\to \subseteq \mathbf{Config} \times \mathbf{Config}$ (and its
reflexive and transitive closure $\to^*$). In this lecture we will explore an alternative
approach---\emph{large-step} operational semantics---which yields the final result of
evaluating an expression directly.

Defining a large-step semantics boils down to specifying a relation $\Downarrow$ that captures
the evaluation of an expression. The $\Downarrow$ relation has the following type:
\[
\Downarrow \; \subseteq \; (\mathbf{Store} \times \mathbf{Exp}) \times (\mathbf{Store} \times
\mathbf{Int}).
\]
We write $\langle \sigma, e \rangle \Downarrow \langle \sigma', n \rangle$ to indicate that
$((\sigma, e), (\sigma', n)) \in {\Downarrow}$. In other words, the expression $e$ with store
$\sigma$ evaluates in one big step to the final store $\sigma'$ and integer $n$.

We define the relation $\Downarrow$ inductively, using inference rules:
\[
\frac{}{\langle \sigma, n \rangle \Downarrow \langle \sigma, n \rangle} \; \text{INT}
\qquad
\frac{n = \sigma(x)}{\langle \sigma, x \rangle \Downarrow \langle \sigma, n \rangle} \;
\text{VAR}
\]
\nt{
Note that in the big step semantics, we need an inference rule for every valid configuration, hence the \textsc{INT} rule.
}
\[
\frac{\langle \sigma, e_1 \rangle \Downarrow \langle \sigma', n_1 \rangle \qquad \langle
\sigma', e_2 \rangle \Downarrow \langle \sigma'', n_2 \rangle \qquad n = n_1 + n_2}{\langle
\sigma, e_1 + e_2 \rangle \Downarrow \langle \sigma'', n \rangle} \; \text{ADD}
\]
\[
\frac{\langle \sigma, e_1 \rangle \Downarrow \langle \sigma', n_1 \rangle \qquad \langle
\sigma', e_2 \rangle \Downarrow \langle \sigma'', n_2 \rangle \qquad n = n_1 \times n_2}{\langle
\sigma, e_1 * e_2 \rangle \Downarrow \langle \sigma'', n \rangle} \; \text{MUL}
\]
\[
\frac{\langle \sigma, e_1 \rangle \Downarrow \langle \sigma', n_1 \rangle \qquad \langle
\sigma'[x \mapsto n_1], e_2 \rangle \Downarrow \langle \sigma'', n_2 \rangle}{\langle \sigma, x
:= e_1 ; e_2 \rangle \Downarrow \langle \sigma'', n_2 \rangle} \; \text{ASSGN}
\]

\ex{Big-step proof}{
To illustrate the use of these rules, consider the following proof tree, which shows that
evaluating $\langle \sigma, \mathit{foo} := 3 ; \mathit{foo} * \mathit{bar} \rangle$ using a
store $\sigma$ such that $\sigma(\mathit{bar}) = 7$ yields $\sigma' = \sigma[\mathit{foo}
\mapsto 3]$ and $21$ as a result:
\[
\cfrac{
\cfrac{}{\langle \sigma, 3 \rangle \Downarrow \langle \sigma, 3 \rangle} \; \text{INT}
\qquad
\cfrac{
\cfrac{}{\langle \sigma', \mathit{foo} \rangle \Downarrow \langle \sigma', 3 \rangle} \;
\text{VAR}
\qquad
\cfrac{}{\langle \sigma', \mathit{bar} \rangle \Downarrow \langle \sigma', 7 \rangle} \;
\text{VAR}
}{\langle \sigma', \mathit{foo} * \mathit{bar} \rangle \Downarrow \langle \sigma', 21
\rangle} \; \text{MUL}
}{\langle \sigma, \mathit{foo} := 3 ; \mathit{foo} * \mathit{bar} \rangle \Downarrow \langle
\sigma', 21 \rangle} \; \text{ASSGN}
\]

A closer look to this structure reveals the relation between small step and large-step
evaluation: a depth-first traversal of the large-step proof tree yields the sequence of
one-step transitions in small-step evaluation.}

\nt{
  Big step semantics don't capture the computational intricacys (if not well-formed it's undefined instead of blocking), and it strugles with concurrency. For this reason it can be used toghether with small step for the critical parts.
}

\nt{
  Also, small step semantics can be used to precysely count computation ticks (cslib 4 lean) for tighter complexity bounds.
}

\section{Equivalence of semantics}

A natural question to ask is whether the small-step and large-step semantics are equivalent.
The next theorem answers this question affirmatively.

\thm{Equivalence of semantics}{
For all expressions $e$, stores $\sigma$ and $\sigma'$, and integers $n$ we have:
\[
\langle \sigma, e \rangle \Downarrow \langle \sigma', n \rangle \quad \text{if and only if}
\quad \langle \sigma, e \rangle \to^* \langle \sigma', n \rangle
\]
}


To streamline the proof, we will work with the following definition of the multi-step
relation:
\[
\frac{}{\langle \sigma, e \rangle \to^* \langle \sigma, e \rangle} \; \text{REFL}
\qquad
\frac{\langle \sigma, e \rangle \to \langle \sigma', e' \rangle \qquad \langle \sigma', e'
\rangle \to^* \langle \sigma'', e'' \rangle}{\langle \sigma, e \rangle \to^* \langle \sigma'',
e'' \rangle} \; \text{TRANS}
\]

\pf{Proof Sketch}{
We show each direction separately.

\medskip
\noindent$\boldsymbol{\implies}$\textbf{:} We want to prove that the following property $P$
holds for all expressions $e \in \mathbf{Exp}$:
\[
P(e) \triangleq \forall \sigma, \sigma' \in \mathbf{Store}.\ \forall n \in \mathbf{Int}.\
\langle \sigma, e \rangle \Downarrow \langle \sigma', n \rangle \implies \langle \sigma, e
\rangle \to^* \langle \sigma', n \rangle
\]
We proceed by structural induction on $e$. We have to consider each of the possible axioms and
inference rules for constructing an expression.

\medskip
\noindent\textbf{Case $e = x$:} Assume that $\langle \sigma, x \rangle \Downarrow \langle
\sigma', n \rangle$. That is, there is some derivation in the large-step operational semantics
whose conclusion is $\langle \sigma, x \rangle \Downarrow \langle \sigma, n \rangle$. There is
only one rule whose conclusion matches the configuration $\langle \sigma, x \rangle$: the
large-step rule \textsc{Var}. Thus, we have $n = \sigma(x)$ and $\sigma' = \sigma$. By the
small-step rule \textsc{Var}, we also have $\langle \sigma, x \rangle \to \langle \sigma, n
\rangle$. By the \textsc{Refl} and \textsc{Trans} rules, we conclude that $\langle \sigma, x
\rangle \to^* \langle \sigma, n \rangle$, which finishes the case.

\medskip
\noindent\textbf{Case $e = n$:} Assume that $\langle \sigma, n \rangle \Downarrow \langle
\sigma', n' \rangle$. There is only one rule whose conclusion matches $\langle \sigma, n
\rangle$: the large-step rule \textsc{Int}. Thus, we have $n' = n$ and $\sigma' = \sigma$ and so
$\langle \sigma, n \rangle \to^* \langle \sigma, n \rangle$ by the \textsc{Refl} rule.

\medskip
\noindent\textbf{Case $e = e_1 + e_2$:} This is an inductive case. We want to prove that if
$P(e_1)$ and $P(e_2)$ hold, then $P(e)$ also holds. Let's write out $P(e_1)$, $P(e_2)$, and
$P(e)$ explicitly.
\[
P(e_1) = \forall n, \sigma, \sigma'.\ \langle \sigma, e_1 \rangle \Downarrow \langle \sigma', n
\rangle \implies \langle \sigma, e_1 \rangle \to^* \langle \sigma', n \rangle
\]
\[
P(e_2) = \forall n, \sigma, \sigma'.\ \langle \sigma, e_2 \rangle \Downarrow \langle \sigma', n
\rangle \implies \langle \sigma, e_2 \rangle \to^* \langle \sigma', n \rangle
\]
\[
P(e) = \forall n, \sigma, \sigma'.\ \langle \sigma, e_1 + e_2 \rangle \Downarrow \langle
\sigma', n \rangle \implies \langle \sigma, e_1 + e_2 \rangle \to^* \langle \sigma', n \rangle
\]
Assume that $P(e_1)$ and $P(e_2)$ hold. Also assume that there exist $\sigma$, $\sigma'$ and
$n$ such that $\langle \sigma, e_1 + e_2 \rangle \Downarrow \langle \sigma', n \rangle$. We need
to show that $\langle \sigma, e_1 + e_2 \rangle \to^* \langle \sigma', n \rangle$.

We assumed that $\langle \sigma, e_1 + e_2 \rangle \Downarrow \langle \sigma', n \rangle$. This
means that there is some derivation whose conclusion is $\langle \sigma, e_1 + e_2 \rangle
\Downarrow \langle \sigma', n \rangle$. By inspection, we see that only one rule has a
conclusion of this form: the \textsc{Add} rule. Thus, the last rule used in the derivation was
\textsc{Add} and it must be the case that $\langle \sigma, e_1 \rangle \Downarrow \langle
\sigma'', n_1 \rangle$ and $\langle \sigma'', e_2 \rangle \Downarrow \langle \sigma', n_2
\rangle$ hold for some $n_1$ and $n_2$ with $n = n_1 + n_2$.

By the induction hypothesis $P(e_1)$, as $\langle \sigma, e_1 \rangle \Downarrow \langle
\sigma'', n_1 \rangle$, we must have $\langle \sigma, e_1 \rangle \to^* \langle \sigma'', n_1
\rangle$. Likewise, by induction hypothesis $P(e_2)$, we have $\langle \sigma'', e_2 \rangle
\to^* \langle \sigma', n_2 \rangle$. By Lemma~1 below, we have,
\[
\langle \sigma, e_1 + e_2 \rangle \to^* \langle \sigma'', n_1 + e_2 \rangle,
\]
and by another application of Lemma~1 we have:
\[
\langle \sigma'', n_1 + e_2 \rangle \to^* \langle \sigma', n_1 + n_2 \rangle
\]
Then, using the small-step \textsc{Add} rule and the multi-step \textsc{Trans} rule, we have:
\[
\cfrac{
\cfrac{n = n_1 + n_2}{\langle \sigma', n_1 + n_2 \rangle \to \langle \sigma', n \rangle} \;
\text{ADD}
\qquad
\cfrac{}{\langle \sigma', n \rangle \to^* \langle \sigma', n \rangle} \; \text{REFL}
}{\langle \sigma', n_1 + n_2 \rangle \to^* \langle \sigma', n \rangle} \; \text{TRANS}
\]
Finally, by two applications of Lemma~2, we obtain $\langle \sigma, e_1 + e_2 \rangle \to^*
\langle \sigma', n \rangle$, which finishes the case.

\medskip
\noindent\textbf{Case $e = e_1 * e_2$.} Similar to case for $e_1 + e_2$ above.

\medskip
\noindent\textbf{Case $e = x := e_1; e_2$.} Omitted. Try it as an exercise.

\medskip
\noindent$\boldsymbol{\impliedby}$\textbf{:} We proceed by induction on the derivation of
$\langle \sigma, e \rangle \to^* \langle \sigma', n \rangle$ with a case analysis on the last
rule used.

\medskip
\noindent\textbf{Case \textsc{Refl}:} Then $e = n$ and $\sigma' = \sigma$. We immediately have
$\langle \sigma, n \rangle \Downarrow \langle \sigma, n \rangle$ by the large-step rule
\textsc{Int}.

\medskip
\noindent\textbf{Case \textsc{Trans}:} Then $\langle \sigma, e \rangle \to \langle \sigma'',
e'' \rangle$ and $\langle \sigma'', e'' \rangle \to^* \langle \sigma', n \rangle$. In this
case, the induction hypothesis gives $\langle \sigma'', e'' \rangle \Downarrow \langle \sigma',
n \rangle$. The result follows from Lemma~3 below.
}

\mlenma{}{
If $\langle \sigma, e \rangle \to^* \langle \sigma', n \rangle$, then the following hold:
\begin{itemize}
    \item $\langle \sigma, e + e_2 \rangle \to^* \langle \sigma', n + e_2 \rangle$
    \item $\langle \sigma, e * e_2 \rangle \to^* \langle \sigma', n * e_2 \rangle$
    \item $\langle \sigma, n_1 + e \rangle \to^* \langle \sigma', n_1 + n \rangle$
    \item $\langle \sigma, n_1 * e \rangle \to^* \langle \sigma', n_1 * n \rangle$
    \item $\langle \sigma, x := e ; e_2 \rangle \to^* \langle \sigma', x := n ; e_2 \rangle$
\end{itemize}
}
\pf{}{
Omitted; try it as an exercise.
}

\mlenma{}{
If $\langle \sigma, e \rangle \to^* \langle \sigma', e' \rangle$ and $\langle \sigma', e'
\rangle \to^* \langle \sigma'', e'' \rangle$, then $\langle \sigma, e \rangle \to^* \langle
\sigma'', e'' \rangle$.
}
\mlenma{}{
Omitted; try it as an exercise.
}

\mlenma{}{
If $\langle \sigma, e \rangle \to \langle \sigma'', e'' \rangle$ and $\langle \sigma'', e''
\rangle \Downarrow \langle \sigma', n \rangle$, then $\langle \sigma, e \rangle \Downarrow
\langle \sigma', n \rangle$.
}
\begin{proof}
Omitted; try it as an exercise.
\end{proof}

% \end{document}
